% page 224 du fac-simile (image p-223 du scan)
% bloc 160.0 x 233.3 mm ; marge gauche 8.0 mm
\pgc{-12.700mm}{0.000mm}{
\l{\cel{3}216.}
\l{}
\l{\sou{hidromekaniko}. Cienco pri la aparati en qui aquo uzesas kam}
\l{\cel{3}forco. - DEFIS.}
\l{}
\l{\sou{hidrom}etro. Instrumento per qua onu mezuras la dikeso di la}
\l{\cel{3}aquo-strato qua pluve falis dum un yaro en ta o ca loko.-}
\l{\cel{3}DEFIRS.}
\l{}
\l{hidrometrio. La cienco qua traktas la mezuro di la denseso e}
\l{\cel{3}di la fluo-rapideso di la liquidi. - DEFIRS.}
\l{}
\l{\sou{hidropatio}. Sistemo terapiala di la mediki qui probas kuracar}
\l{\cel{3}nur per la uzo (sur ed en la korpo) di aquo pura. - DEFIRS.}
\l{}
\l{\sou{hidropso}. (\sou{patol.}) Akumulajo morbala ek liquido seroza en\cel{2}ka-}
\l{\cel{3}vajo di la korpo, od en la tisuo celulara. - DEFIS.}
\l{}
\l{\sou{hidroquinono}. (\sou{kemio}). C6 H4 (OH)2.}
\l{}
\l{\sou{hidroskopo}. La persono qua profesione okupesas per la arto}
\l{\cel{3}deskovrar la aquo-fonti. - DEFIRS.}
\l{}
\l{\sou{hidroskopio}. La arto deskovrar la aquo-fonti. - DFIRS.}
\l{}
\l{\sou{hidrostatiko}. La parto di hidrauliko qua traktas la equilibro}
\l{\cel{3}di la gasi e di la liquidi, e la presi da li sur la parie-}
\l{\cel{3}ti di la recipienti qui kontenas li. - DEFIRS.}
\l{}
\l{\sou{hidroterapio}. (\sou{medic.}) Kuraco di ula morbi per apliko di aquo}
\l{\cel{3}sur la korpo (balni, dushi, vapor-ago, e c.) - DEFIRS.}
\l{}
\l{\sou{hieno}. (\sou{zool.}) Mamifero karnavida, digitigrada, analoga a}
\l{\cel{3}volfo. - L.\cel{1}\sou{hyaena}. - DEFIRS.}
\l{}
\l{\sou{hiero}. La dio qua preiras nemediate olta en qua onu vivas. FIS.}
\l{}
\l{\sou{hierarkio}. I. (\sou{kristanismo)}.\cel{2}Ordino di la subordineso di la}
\l{\cel{3}ciel-kori, di la anjeli. - II. Ordino di la subordineso di}
\l{\cel{3}ti qui esas situita en rangi ne-interegala. - DEFIRS.}
\l{}
\l{\sou{hieratika}.\cel{2}Qua koncernas la kozi sakra. - (\sou{che la Egiptiani}}
\l{\cel{3}\sou{antiqua}). (Skribo-stilo) trasebla rapide, por la linguo}
\l{\cel{3}hieroglifala. - DEFIRS.}
\l{}
\l{\sou{hieroglifo}. I. (\sou{arkeol.}) En la skribo-stilo di la Egiptiani an-}
\l{\cel{3}tiqua, signo sakra qua expresis simbole la idei, per la re-}
\l{\cel{3}prezento di ula objekti naturala. - II. (\sou{metaf}.)Kozo enig-}
\l{\cel{3}mata. - DEFIRS.}
\l{}
\l{\textquotedbl{}high\cel{1}\sou{lif}e\textquotedbl{}. Granda mondumo, ek la aristokrati di la socio.}
\l{}
\l{\sou{higieno}.\cel{2}Parto de medicino, qua traktas la rejimo adoptenda}
\l{\cel{3}da la homo por konservar lua saneso. - DEFIRS.}
}
